% Irodov user-provided PDF pp.69--70, original printed pp.60--61.
% Relativity section 1.8, original problems 371--380.
\exproblem{371}{60}
质量为 $m_0$ 的粒子以何速率运动时，
其相对论动量模 $p$ 是牛顿动量模
$m_0v$ 的 $\eta=2.0$ 倍？
\exsolution
由相对论动量
$p=\gamma m_0v$，
其中 $\gamma=(1-\beta^2)^{-1/2}$、
$\beta=v/c$，
有
\[
\frac{p}{m_0v}=\gamma=\eta.
\]
因此
\[
\boxed{v=c\sqrt{1-\eta^{-2}}
=\frac c\eta\sqrt{\eta^2-1}
=\frac{\sqrt3}{2}c\simeq0.866c}.
\]
\excheck $p/(m_0v)$ 随速率增长且始终不小于1；
题给 $\eta=2$ 对应 $\gamma=2$，
与第373题的动能条件巧合相同。

\exproblem{372}{60}
把静质量为 $m_0$ 的粒子从
$v_1=0.60c$ 加速到 $v_2=0.80c$，
外力应做多少功？将相对论结果与
牛顿动能增量比较。
\exsolution
外力功等于总相对论能量的增量：
\[
W=(\gamma_2-\gamma_1)m_0c^2,
\quad
\gamma_i=\frac1{\sqrt{1-v_i^2/c^2}}.
\]
$\gamma_1=1/\sqrt{1-0.60^2}=5/4$，
$\gamma_2=1/\sqrt{1-0.80^2}=5/3$，
所以
\[
\boxed{W=\left(\frac53-\frac54\right)m_0c^2
=\frac5{12}m_0c^2\simeq0.42m_0c^2}.
\]
若误用牛顿动能，得到
\[
W_{\rm Newt}=\frac12m_0(v_2^2-v_1^2)
=\tfrac12m_0c^2(0.64-0.36)
=\boxed{0.14m_0c^2}.
\]
\excheck 牛顿估计仅约为正确值的三分之一，
说明在 $0.6c$ 到 $0.8c$ 的速度范围
不能忽略相对论效应。

\exproblem{373}{60}
粒子速率达到多少时，其动能
$T$ 恰好等于其静止能量 $E_0=m_0c^2$？
\exsolution
粒子的动能为
$T=(\gamma-1)m_0c^2$。
令 $T=m_0c^2$，得 $\gamma=2$。
由
$\gamma^{-2}=1-\beta^2$，
\[
\boxed{v=c\sqrt{1-\frac1{\gamma^2}}
=\frac{\sqrt3}{2}c
\simeq2.60\times10^8\,\mathrm{m/s}}.
\]
\excheck 动能等于静能并不表示粒子达到光速；
有静质量粒子的速率始终小于 $c$。

\exproblem{374}{60}
当粒子动能与静止能量之比
$q=T/(m_0c^2)$ 足够小时，
能否用牛顿公式 $T=m_0v^2/2$ 估算其速度？
求使速率相对误差小于
$\varepsilon=0.010$ 的 $q$ 上界近似值。
\exsolution
精确相对论关系
$T=(\gamma-1)m_0c^2$
给出 $\gamma=1+q$，从而
\[
v=c\sqrt{1-(1+q)^{-2}}.
\]
由牛顿动能反算的速率为
$v_{\rm N}=c\sqrt{2q}$。
对 $q\ll1$，
\[
v=c\sqrt{2q}
\left(1-\frac34q+O(q^2)\right).
\]
于是采用真实速度作相对误差分母，
\[
\frac{v_{\rm N}-v}{v}
=\frac34q+O(q^2).
\]
若要求误差小于 $\varepsilon$，
低能极限需要
\[
\boxed{\frac{T}{m_0c^2}=q
\lesssim\frac43\varepsilon
=\frac43(0.010)
\simeq0.013}.
\]
\excheck 这是小 $q$ 的\emph{一阶误差界}，
并非忽略高阶项后的严格解析不等式。
如需严密阈值，应将
$v_{\rm N}/v-1<\varepsilon$
直接代入精确公式求解。

\exproblem{375}{60}
推导静质量为 $m_0$ 的粒子的
相对论动量 $p$ 与动能 $T$ 的关系。
并计算动能为 $500\,\mathrm{MeV}$
的质子的动量模。
\exsolution
相对论总能量
$E=T+m_0c^2$，
满足
$E^2=p^2c^2+m_0^2c^4$。
故
\[
p^2c^2=(T+m_0c^2)^2-m_0^2c^4
=T(T+2m_0c^2),
\]
因此
\[
\boxed{p=\frac1c
\sqrt{T(T+2m_0c^2)}}.
\]
质子静止能量
$m_pc^2\simeq938\,\mathrm{MeV}$，
代入 $T=500\,\mathrm{MeV}$：
\[
pc=\sqrt{500(500+2\cdot938)}\,\mathrm{MeV}
\simeq1090\,\mathrm{MeV},
\]
\[
\boxed{p\simeq1.09\,\mathrm{GeV}/c}.
\]
\excheck 若使用牛顿公式 $p=\sqrt{2m_0T}$
会漏掉 $T^2$ 项，无法在高动能下保证精度。

\exproblem{376}{60}
一束动能为 $T$ 的相对论性带电粒子
垂直射向完全吸收靶，
束流电流为 $I$，
每粒子电荷大小为 $e$、静质量为 $m_0$。
求靶受到的平均冲击力，以及靶释放的功率。
若给定受照面积 $A$，求平均压强。
\exsolution
单位时间打到靶上的粒子数为
$\dot N=I/e$。
粒子被完全吸收时，其初始动量全部传给靶，
每粒子动量模由能量—动量关系为
\[
p=\frac1c\sqrt{T(T+2m_0c^2)}.
\]
平均冲击力是动量通量：
\[
\boxed{F=\frac{I}{ec}
\sqrt{T(T+2m_0c^2)}}.
\]
若束斑面积为 $A$，
压力（压强）为
$\boxed{P_{\rm press}=F/A}$；
\emph{原题未提供面积，不能单独算出数值压强}。
假定打到靶上的粒子仅把其\emph{动能}
转化为热，释放的功率
\[
\boxed{\mathcal P=\dot NT=\frac{IT}{e}}.
\]
\excheck 原书简答使用符号 $F$ 表示作用力，
而不是面积归一化压强。若讨论粒子静质量
经反应转换为热，还需给出具体反应过程，
本题只考虑吸收、停止。

\exproblem{377}{60}
一球以相对论速率 $v$ 在稀薄气体中运动，
气体静止系每单位体积含 $n$ 个慢速粒子，
各粒子静质量为 $m$。
假设气体粒子在球表面垂直于球速的微面元上
发生理想弹性镜面反射。
求该面元受到的气体压强 $p$，
并验证在球的瞬时静止系和气体静止系所得相同。
\exsolution
记 $\gamma=(1-v^2/c^2)^{-1/2}$。
先到\emph{球的静止系}：
气体朝向表面的速度为 $v$，
长度收缩使粒子数密度变为
$n'=\gamma n$。
每粒子入射动量模为
$p_i=\gamma mv$，弹性反射使法向动量反号，
故每次冲量为 $2\gamma mv$。
单位面积单位时间碰撞数为
$n'v=\gamma nv$，因此
\[
\boxed{p=(\gamma nv)(2\gamma mv)
=\frac{2nmv^2}{1-v^2/c^2}}.
\]
在\emph{气体静止系}中，靶面以速度 $v$
迎向静止粒子。反射后粒子的速率大小由
相对论速度合成得
$u=2v/(1+v^2/c^2)$，
其动量为
$\gamma_u m u=2\gamma^2mv$。
单位面积单位时间被迎面扫过的粒子数为 $nv$，
所以靶面所得冲量通量仍为
\[
(nv)(2\gamma^2mv)=2\gamma^2nmv^2.
\]
\excheck 这里是法向的\emph{压力即动量流率密度}；
在选定的理想单向弹性反射模型下两惯性系计算一致。
粒子不能被当作“用牛顿动量 $mu$”来计算，
否则少一个相对论因子。

\exproblem{378}{60}
静质量为 $m_0$ 的粒子从
$t=0$ 时刻由静止开始受到
沿直线方向大小恒定的力 $F$。
求粒子速率 $v(t)$ 及在力方向的位移 $x(t)$。
\exsolution
沿运动方向的相对论动量为
$p=\gamma m_0v$，
动力学方程 $\mathrm dp/\mathrm dt=F$。
由 $p(0)=0$ 积分：
$p(t)=Ft$。
又
$v=pc^2/\sqrt{m_0^2c^4+p^2c^2}$，
因此
\[
\boxed{v(t)=
\frac{F t}{\sqrt{m_0^2+(F t/c)^2}}}.
\]
从 $x(0)=0$ 出发，
\[
x(t)=\int_0^t\frac{F s\,\mathrm ds}
{\sqrt{m_0^2+(Fs/c)^2}}
=\frac{m_0c^2}{F}
\left[\sqrt{1+\left(\frac{F t}{m_0c}\right)^2}-1\right].
\]
故
\[
\boxed{x(t)=\frac{
\sqrt{m_0^2c^4+F^2c^2t^2}-m_0c^2}{F}}.
\]
\excheck 在 $Ft\ll m_0c$ 时，
$v\simeq(F/m_0)t$，
$x\simeq Ft^2/(2m_0)$；
而 $t\to\infty$ 时
$v\to c^-$，始终不超光速。

\exproblem{379}{60}
一个粒子在 $K$ 系的运动规律为
$x(t)=\sqrt{a^2+c^2t^2}$，
其中 $a>0$ 为常长度、
$c$ 是光速。
求作用于它的相对论力 $F$。
\exsolution
直接微分：
\[
v=\dot x=\frac{c^2t}{x},\qquad
w=\ddot x=\frac{c^2a^2}{x^3}.
\]
由
$1-v^2/c^2=1-c^2t^2/x^2=a^2/x^2$
得洛伦兹因子
$\gamma=x/a$。
沿速度方向的动量微分
\[
F=\frac{\mathrm d}{\mathrm dt}(\gamma m_0v)
=m_0\gamma^3 w
=m_0\left(\frac xa\right)^3
\frac{c^2a^2}{x^3}.
\]
故作用力为恒定值
\[
\boxed{F=\frac{m_0c^2}{a}}
\]
且沿 $+x$ 方向。
\excheck 虽然坐标加速度 $w(t)$ 不恒定，
相对论力 $F=\mathrm dp/\mathrm dt$ 却恒定；
不能将牛顿式 $F=m_0w$ 直接使用。

\exproblem{380}{60--61}
从相对论动力学基本方程出发：
(a) 说明哪些情况下粒子的加速度方向
与作用力 $\bm F$ 相同；
(b) 当 $\bm F\perp\bm v$ 与
$\bm F\parallel\bm v$ 时，
求力与加速度模之间的比例因子。
\exsolution
相对论动量
$\bm p=\gamma m_0\bm v$。
求导并使用
$\dot\gamma=\gamma^3(\bm v\cdot\bm w)/c^2$，
得到
\[
\bm F
=\frac{\mathrm d\bm p}{\mathrm dt}
=\gamma m_0\bm w+
\frac{\gamma^3m_0}{c^2}
(\bm v\cdot\bm w)\bm v.
\]
将加速度分解为沿速率方向和法向分量
$\bm w=\bm w_\parallel+\bm w_\perp$，
可化为
\[
\bm F_\perp=\gamma m_0\bm w_\perp,\qquad
\bm F_\parallel=\gamma^3m_0\bm w_\parallel.
\]
由于两个系数通常不同，
$\bm F$ 与 $\bm w$ 在一般斜向情形并不同向。
仅在 $\bm F\parallel\bm v$
或 $\bm F\perp\bm v$
（以及 $v=0$ 的退化瞬间）时，二者同向。
因此
\[
\boxed{
F_\perp=\frac{m_0w_\perp}{\sqrt{1-\beta^2}},
\qquad
F_\parallel=\frac{m_0w_\parallel}{(1-\beta^2)^{3/2}}
},\qquad \beta=v/c.
\]
\excheck 纵向“等效惯性系数”为 $\gamma^3m_0$，
横向为 $\gamma m_0$；这些是方向依赖的动力学关系，
不是两个不同的不变静质量。
